\chapter{绪论}

随着汽车工业的发展，汽车在给人们带来极大便利的同时,汽车在大量消耗石油资源的同时,尾气排放含有有害还严重的污染，已成为城市大气的主要污染源,实现资源与环境在可持续发展前提下清洁燃烧已成为当前内燃机研究的重要课题。能源资源是发动机发展的重要基础，也是重要问题。目前，人们对环境污染的关注程度越来越高，因而越来越严格地制定了以NOx和PM为主要污染源的排放法规，例如，欧洲在1998年生效的一项超低排放法规规定：NOx+HC排放<2.5g/bph-hr, PM排放<0.05g/bph-hr。为满足严格的排放要求，研究者们在改善发动机燃烧系统方面进行了大量研究，并取得了一定进展\cite{zhang1998forest}。

传统火花点燃式发动机,尾气排放污染物主要是氮氧化物(NOx)、碳氢化合物(HC)和一氧化碳(CO),虽然通过高效催化转化器可以解决,但要达到欧IV排放标准存在较大难度,而点燃式发动机效率低,具有中高负荷时存在较大的泵气损失,点燃式效率高,但其中的NOx和碳烟微粒排放(PM)却难以控制,使用一种排放较少的措施,往往会增加另一种排放。在不降低发动机热效率, NOx的处理转化技术尚未成熟。点燃式和压燃式燃烧方式往往存在碳烟和氮氧化物生成的trade-off关系，因此需要探索新的燃烧模式，通过改进燃烧同时降低碳烟和氮氧化物排放。

\section{HCCI燃烧数值模拟研究现状}

HCCI燃烧着火和燃烧放热过程与传统的火花点燃式和压缩点燃式发动机有着本质区别。在HCCI燃烧着火和燃烧放热过程中，燃料的化学反应和化学动力学起着主要作用。因此，在传统燃烧的数值模拟中主要采用的零维和准维燃烧模型对于HCCI燃烧着火模型的建立需要主要依靠燃料的反应机理和化学动力学模型。

\subsection{HCCI数值模拟模型}

目前HCCI数值模拟研究主要集中在单区模型和多维模型等\cite{chineseacademy1990}，现将各种模型直接优缺点介绍如下：

\begin{enumerate}
	\item 单区模型

	单区模型将整个气缸视为一个均匀的热力学系统，假设缸内混合气在压缩、燃烧和膨胀过程中始终处于热力学平衡状态。该模型计算简单，能够快速预测着火时刻和燃烧放热率，但无法考虑缸内温度和浓度分布的不均匀性。

	\item 双区多区模型

	双区模型将缸内划分为已燃区和未燃区，多区模型则进一步细分为多个区域，每个区域具有不同的温度和浓度。该模型在计算效率和精度之间取得了较好的平衡，能够部分考虑缸内不均匀性的影响。

	\item 多维模型

	多维模型（包括二维和三维CFD模型）通过求解流体力学控制方程，能够详细描述缸内流场、温度场和浓度场的空间分布。该模型计算精度高，但计算量巨大，通常需要与化学动力学模型耦合使用。
\end{enumerate}

\section{本文研究内容}

本文主要研究DME均质压燃燃烧的数值模拟方法，具体研究内容包括：

\begin{enumerate}
	\item 建立DME均质压燃燃烧的单区数值模型；
	\item 研究压缩比、当量比、进气温度等参数对HCCI着火和燃烧的影响；
	\item 研究EGR和燃料添加剂对着火时刻的影响规律。
\end{enumerate}
